\documentclass[UTF8]{ctexart}
\usepackage{xeCJK}
\usepackage{geometry}
\geometry{left=2cm,right=2cm,top=2.5cm,bottom=2.5cm}
\setlength{\headheight}{12.7pt}

\usepackage{titlesec}
\titleformat{\section}
  {\normalfont\large\bfseries}
  {\thesection}
  {1em}
  {}
\titlespacing*{\section}
  {0pt}
  {3.5ex plus 1ex minus .2ex}
  {2.3ex plus .2ex}

\usepackage{amsmath}
\usepackage{amssymb}
\usepackage{amsthm}
\usepackage{enumitem}
\setlist[enumerate]{itemsep=0pt,topsep=0pt,parsep=0pt,leftmargin=0pt,itemindent=2.5em}
\setlist[itemize]{itemsep=0pt,topsep=0pt,parsep=0pt,leftmargin=0pt,itemindent=2.5em}
\usepackage{fancyhdr}
\pagestyle{fancy}
\fancyhf{}
\usepackage{graphicx}
\usepackage{hyperref}
\usepackage{indentfirst}
\usepackage{lmodern}


\begin{document}
\setlength{\abovedisplayskip}{1pt}
\setlength{\belowdisplayskip}{1pt}

\section{子集与空集}

\hypertarget{sec:定义1.1}{}
\noindent\textbf{定义1.1 (子集)}

给定集合$a,b$, 如果$\forall x\in a:x\in b$,
那么称$a$是$b$的\textbf{子集}, 记作$a\subset b$.

如果$a\subset b$并且$a\neq b$, 那么称$a$是$b$的\textbf{真子集}, 记作$a\subsetneq b$.

\vspace{10pt}

\hypertarget{sec:定义1.2}{}
\noindent\textbf{定义1.2 (空集)}

如果集合$a$满足$\forall x:x\notin a$, 就称$a$是\textbf{空集}, 记作$\varnothing$.

根据\href{../01 ZFC 公理/01 ZFC 公理#nameddest=sec:分离公理模式}
{分离公理}和\href{../01 ZFC 公理/01 ZFC 公理#nameddest=sec:外延公理}{外延公理},
空集是存在且唯一的.





\vspace{20pt}

\section{对集}

\hypertarget{sec:定义2.1}{}
\noindent\textbf{定义2.1 (对集)}

给定集合$a,b$, 如果$c$满足$\forall x:(x\in c\leftrightarrow x=a\lor x=b)$,
就称$c$是$a,b$的\textbf{对集}, 记作$\{a,b\}$.

根据\href{../01 ZFC 公理/01 ZFC 公理#nameddest=sec:对集公理}{对集公理}和
外延公理, 任意$a,b$的对集是存在且唯一的.

\vspace{10pt}

显然$\{a,b\}=\{b,a\}$. 把$\{a,a\}$简记为$\{a\}$, 称之为$a$的\textbf{单点集}.

\vspace{10pt}

\hypertarget{sec:定理2.1}{}
\noindent\textbf{定理2.1}

对任意的集合$a,b,b'$, 如果$\{a,b\}=\{a,b'\}$, 那么$b=b'$.

\noindent{\kaishu 证明.}

\begin{enumerate}
    \item 如果$a=b$, 那么$b'\in\{a,b\}=\{a\}$, 即$b'=a=b$;
    \item 如果$a\neq b$, 那么$b'\in\{a,b\}$可得$b'=a\lor b'=b$, 从而也有$b'=b$.
    \hfill $\qedsymbol$
\end{enumerate}

\vspace{10pt}

\hypertarget{sec:定义2.2}{}
\noindent\textbf{定义2.2 (有序对)}

给定集合$a,b$, 定义它们的\textbf{有序对}为
\[
    (a,b)\overset{\mathrm{def}}{=}\{\{a\},\{a,b\}\}.
\]

\vspace{10pt}

有序对可以唯一确定两个原集合.

\vspace{10pt}

\hypertarget{sec:定理2.2}{}
\noindent\textbf{定理2.2}

对任意的集合$a,b,a',b'$, 如果$(a,b)=(a',b')$, 那么$a=a'$, $b=b'$.

\noindent{\kaishu 证明.}

依定义, $\{a'\}\in\{\{a'\},\{a',b'\}\}=\{\{a\},\{a,b\}\}$.

\begin{enumerate}
    \item 如果$\{a'\}=\{a\}$, 那么$\{a',b'\}=\{a,b\}$, 从而$a'=a$, $b'=b$;
    \item 如果$\{a'\}=\{a,b\}$, 那么$\{a',b'\}=\{a\}$, 从而$b=a'=a=b'$.
    \hfill $\qedsymbol$
\end{enumerate}

\vspace{10pt}

由此定理, 对任意的有序对$p$, 存在唯一的$x,y$使得$p=(x,y)$.





\newpage

\section{分离与替换}

给定集合$a,p_1,\cdots,p_n$.

如果$P$是扩充$n+1$元关系符号, 那么根据分离公理和外延公理, 存在唯一的集合$b$使得
\[
    \forall x:(x\in b\leftrightarrow(x\in a\land P(p_1,\cdots,p_n,x))),
\]
记为$\{x\in a\mid P(p_1,\cdots,p_n,x)\}$.

如果$F$是扩充$n+1$元函数符号, 那么
根据\href{../01 ZFC 公理/01 ZFC 公理#nameddest=sec:替换公理模式}{替换公理}和
外延公理, 存在唯一的集合$b$使得
\[
    \forall y:(y\in b\leftrightarrow\exists x\in a:y=F(p_1,\cdots,p_n,x)),
\]
记为$\{F(p_1,\cdots,p_n,x)\mid x\in a\}$.





\vspace{20pt}

\section{并集与幂集}

\hypertarget{sec:定义4.1}{}
\noindent\textbf{定义4.1 (并集)}

给定集合$a$, 如果$b$满足$\forall x:(x\in b\leftrightarrow\exists y\in a:x\in y)$,
就称$b$是$a$的\textbf{并集}, 记作$\displaystyle\bigcup a$.

根据\href{../01 ZFC 公理/01 ZFC 公理#nameddest=sec:并集公理}{并集公理}和
外延公理, 任意$a$的并集是存在且唯一的.

\vspace{10pt}

\hypertarget{sec:定义4.2}{}
\noindent\textbf{定义4.2 (交集)}

给定集合$a$, 定义它的\textbf{交集}
\[
    \bigcap a\overset{\mathrm{def}}{=}\left\{x\in\bigcup a\,\middle\vert\,
    \forall y\in a:x\in y\right\}.
\]

\vspace{10pt}

显然如果$a$非空, 那么对任意的$x$, $\displaystyle x\in\bigcap a$当且仅当
$\forall y\in a:x\in y$.

\vspace{10pt}

\hypertarget{sec:定义4.3}{}
\noindent\textbf{定义4.3}

给定集合$a,b$, 定义:
\begin{enumerate}
    \item $\displaystyle a\cup b\overset{\mathrm{def}}{=}\bigcup\{a,b\}$,
    显然对任意的$x$, $x\in a\cup b$当且仅当$x\in a\lor x\in b$,
    \item $\displaystyle a\cap b\overset{\mathrm{def}}{=}\bigcap\{a,b\}$,
    显然对任意的$x$, $x\in a\cap b$当且仅当$x\in a\land x\in b$,
    \item $\displaystyle a\setminus b\overset{\mathrm{def}}{=}
    \{x\in a\mid x\notin b\}$,
    显然对任意的$x$, $x\in a\setminus b$当且仅当$x\in a\land x\notin b$.
\end{enumerate}

\vspace{10pt}

\hypertarget{sec:定义4.4}{}
\noindent\textbf{定义4.4 (幂集)}

给定集合$a$, 如果集合$b$满足$\forall x:(x\in b\leftrightarrow x\subset a)$,
就称$b$是$a$的\textbf{幂集}, 记作$\mathcal{P}(a)$.

根据\href{../01 ZFC 公理/01 ZFC 公理#nameddest=sec:幂集公理}{幂集公理}和
外延公理, 任意$a$的幂集是存在且唯一的.

\vspace{10pt}

并集与幂集有如下的对偶关系.

\vspace{10pt}

\hypertarget{sec:定理4.1}{}
\noindent\textbf{定理4.1}

对任意的集合$a$, $\displaystyle\bigcup\mathcal{P}(a)=a$,
$\displaystyle\mathcal{P}\left(\bigcup a\right)\supset a$.

\noindent{\kaishu 证明.}

对任意的$x$ :

\begin{enumerate}
    \item $\displaystyle x\in\bigcup\mathcal{P}(a)\iff
    \exists y:(y\in\mathcal{P}(a)\land x\in y)\iff
    \exists y:(y\subset a\land x\in y)\iff x\in a$,
    \item 如果$x\in a$, 显然$\displaystyle x\subset\bigcup a$,
    即$\displaystyle x\in\mathcal{P}\left(\bigcup a\right)$. \hfill $\qedsymbol$
\end{enumerate}





\newpage

\section{二元 Descartes 积}

\hypertarget{sec:定义5.1}{}
\noindent\textbf{定义5.1 (二元Descartes积)}

给定集合$a,b$, 定义它们的\textbf{二元Descartes积}
\[
    a\times b\overset{\mathrm{def}}{=}\{p\in\mathcal{P}\mathcal{P}(a\cup b)
    \mid\exists x\in a\,\exists y\in b:p=(x,y)\}.
\]

\vspace{10pt}

对任意的$x\in a$和$y\in b$, 有
\[\begin{aligned}
    &x,y\in a\cup b\\
    \implies&\{x\},\{x,y\}\subset a\cup b
    \implies\{x\},\{x,y\}\in\mathcal{P}(a\cup b)\\
    \implies&\{\{x\},\{x,y\}\}\subset\mathcal{P}(a\cup b)
    \implies\{\{x\},\{x,y\}\}\in\mathcal{P}\mathcal{P}(a\cup b)\\
    \implies&(x,y)\in\mathcal{P}\mathcal{P}(a\cup b)\\
    \implies&(x,y)\in a\times b,\\[3pt]
\end{aligned}\]
所以对任意的$p$, $p\in a\times b$当且仅当$\exists x\in a\,\exists y\in b:p=(x,y)$.





\vspace{20pt}

\section*{附录1\ \ 多元运算}
\addcontentsline{toc}{section}{(附录1) 多元运算}

为了方便某些情境中的描述, 对集、有序对和二元Descartes积可以扩展到多元的情况.

\vspace{10pt}

\hypertarget{sec:定义A1.1}{}
\noindent\textbf{定义A1.1 (多元集)}

给定集合$a,b,c$, 记$\{a,b,c\}\overset{\mathrm{def}}{=}\{a,b\}\cup\{c\}$.
显然对任意的$x$,
\[
    x\in\{a,b,c\}\leftrightarrow x=a\lor x=b\lor x=c.
\]

给定集合$a,b,c,d$, 记$\{a,b,c,d\}\overset{\mathrm{def}}{=}\{a,b,c\}\cup\{d\}$.
显然对任意的$x$,
\[
    x\in\{a,b,c,d\}\leftrightarrow x=a\lor x=b\lor x=c\lor x=d.
\]

依此类推.

\vspace{10pt}

\hypertarget{sec:定义A1.2}{}
\noindent\textbf{定义A1.2 (有序组)}

给定集合$a,b,c$, 记$(a,b,c)\overset{\mathrm{def}}{=}((a,b),c)$.
显然对任意的$a,b,c,a',b',c'$,
\[
    (a,b,c)=(a',b',c')\to a=a'\land b=b'\land c=c'.
\]

给定集合$a,b,c,d$, 记$(a,b,c,d)\overset{\mathrm{def}}{=}((a,b,c),d)$.
显然对任意的$a,b,c,d,a',b',c',d'$,
\[
    (a,b,c,d)=(a',b',c',d')\to a=a'\land b=b'\land c=c'\land d=d'.
\]

依此类推.

\vspace{10pt}

\hypertarget{sec:定义A1.3}{}
\noindent\textbf{定义A1.3 (多元Descartes积)}

给定集合$a,b,c$, 记$a\times b\times c\overset{\mathrm{def}}{=}
(a\times b)\times c$. 显然对任意的$p$,
\[
    p\in a\times b\times c\leftrightarrow
    \exists x\in a\,\exists y\in b\,\exists z\in c:p=(x,y,z).
\]

给定集合$a,b,c,d$, 记$a\times b\times c\times d\overset{\mathrm{def}}{=}
(a\times b\times c)\times d$. 显然对任意的$p$,
\[
    p\in a\times b\times c\times d\leftrightarrow
    \exists x\in a\,\exists y\in b\,\exists z\in c\,\exists w\in d:p=(x,y,z,w).
\]

依此类推.

\end{document}